% ============================================================
% Parecer da banca Banca_TCC sobre Artigo.tex (RadarTorres)
% Gerado pelo agente banca-tcc (.claude/agents/banca-tcc.md).
% Não editar à mão sem regenerar a partir de uma nova análise.
% Compile com pdfLaTeX.
%
% QUINTA revisão. As anteriores (27/10, nota 6,33; segunda revisão, nota 6,20;
% terceira revisão, nota 7,9; quarta revisão, nota 7,7) estão em
% Parecer_Banca_TCC.tex, Parecer_Banca_TCC_rev3.tex e Parecer_Banca_TCC_rev4.tex
% e permanecem intactas. Esta verifica a reconciliação das branches main e
% banco-dados (achado C.4 da 4ª revisão) e as reescritas de D.2 e D.5.
% ============================================================
\documentclass[11pt,a4paper]{article}

\usepackage[T1]{fontenc}
\usepackage[utf8]{inputenc}
\usepackage[brazil]{babel}
\usepackage[a4paper,margin=2.5cm]{geometry}
\usepackage{xcolor}
\usepackage{longtable}
\usepackage{array}
\usepackage{url}

\title{\textbf{Parecer da Banca Examinadora --- TCC ``RadarTorres''}\\[0.3em]
\large Quinta revisão --- gerado pelo agente \texttt{banca-tcc}}
\author{}
\date{Análise referente a \texttt{Artigo.tex}/\texttt{Artigo.pdf} (RadarTorres, TCC de Engenharia da
Computação, UCL)\\[0.3em]
Auditado contra o código-fonte, os artefatos em JSON e o histórico git, agora com \texttt{main} e
\texttt{banco-dados} reconciliadas num único commit}

\newcommand{\critico}{\textbf{\color{red!60!black}[CRÍTICO]}}
\newcommand{\alto}{\textbf{\color{orange!70!black}[ALTO]}}
\newcommand{\moderado}{\textbf{\color{blue!50!black}[MODERADO]}}
\newcommand{\baixo}{\textbf{\color{gray}[BAIXO]}}
\newcommand{\retirado}{\textbf{[RETIRADO]}}
\newcommand{\resolvido}{\textbf{\color{green!40!black}[RESOLVIDO]}}

\begin{document}
\maketitle

% ============================================================
% QUADRO-RESUMO
% ============================================================
\begin{center}
\fcolorbox{black}{gray!8}{%
\begin{minipage}{0.93\textwidth}
\vspace{0.4em}
\begin{center}
{\Large\textbf{Nota: 8,0 / 10}}\\[0.2em]
{\small pareceres anteriores: 6,33 (1ª) \quad$\cdot$\quad 6,20 (2ª) \quad$\cdot$\quad 7,9 (3ª)
\quad$\cdot$\quad 7,7 (4ª)}
\end{center}
\vspace{0.3em}
\hrule
\vspace{0.6em}
\textbf{Veredicto.} O achado que bloqueava a defesa na 4ª revisão --- bifurcação não reconciliada
entre \texttt{main} e \texttt{banco-dados}, deixando a correção do achado crítico original (D.1) numa
branch que não era a aberta --- está \textbf{resolvido}: as duas branches hoje apontam para o mesmo
commit (\texttt{5ef8f91}), e a frase corrigida sobre a rede não versionada está na branch fisicamente
aberta. D.5 (Tabela de certeza estática) recebeu a resposta correta para um artefato irrecuperável:
declarar o limite na legenda, em vez de inventar um dump. D.2 foi parcialmente corrigido --- a
substituição dos valores pelo parâmetro do experimento está certa, mas a frase nova inverte a ordem
cronológica real entre a elevação dos padrões e a data do experimento, o que é um achado novo, de
gravidade baixa.

\vspace{0.5em}
\textbf{Achados:} 0 críticos abertos \quad$\cdot$\quad 0 altos abertos \quad$\cdot$\quad 1 crítico
resolvido \quad$\cdot$\quad 2 altos resolvidos/rebaixados \quad$\cdot$\quad 3 moderados
\quad$\cdot$\quad 5 baixos \quad$\cdot$\quad 1 retirado
\vspace{0.6em}
\end{minipage}}
\end{center}

\vspace{0.8em}

\textit{Como ler este parecer:} a Seção A resume; a B justifica a nota critério por critério; a C
compara com o parecer anterior; a D lista os problemas, cada um com gravidade no título; a E mostra o
que foi conferido contra o código e os dados, com veredicto por item; a I traz as pendências em
ordem de prioridade; a J diz o que \emph{não} foi verificado. Quem tem pouco tempo: leia A, H e I.

\vspace{0.5em}

\section*{A. Síntese}

Esta rodada foi aberta para verificar três mudanças relatadas desde a 4ª revisão: a reconciliação das
branches \texttt{main} e \texttt{banco-dados}; a reescrita de D.2 (parâmetros de população/espécies);
e a resposta a D.5 (Tabela de certeza estática). As três foram conferidas contra o repositório tal
como está hoje, não contra o relato dos autores.

\textbf{Branches.} \texttt{git merge-base --is-ancestor main banco-dados} e o comando inverso
retornam sucesso nos dois sentidos, e \texttt{git rev-parse main banco-dados} mostra o mesmo hash
(\texttt{5ef8f91c43ed5661e0748e613b2e4d02f3245bf1}) para as duas. Não há mais bifurcação: existe uma
única história, e os commits \texttt{0f58ccb} (remove os genes do versionamento) e \texttt{501ce86}
(reescreve a frase da Seção 2.3.5) estão nela. \texttt{git ls-tree -r HEAD --name-only} não lista mais
\texttt{genoma-controlador.json} nem \texttt{rede-controlador.net}; o \texttt{.gitignore} tem a linha
que os exclui; e o texto em disco, na linha 478--480, diz agora ``a versão entregue não inclui rede
treinada; um executável compilado a partir dela usa a heurística pura, comportamento de fábrica
previsto no projeto'' --- a frase exata que a 3ª e a 4ª revisão recomendaram. O \texttt{.csproj} ganhou
também um comentário explicando a condição \texttt{Exists(...)}, coerente com o que está no disco.
\textbf{D.1 está resolvido.} Os dois arquivos de rede ainda existem fisicamente no disco desta máquina
(resíduo de execuções locais do treinamento), mas não estão rastreados pelo git nem seriam trazidos por
um clone limpo --- o que é exatamente o comportamento que a correção promete.

\textbf{D.2.} A linha 235 hoje diz: ``parâmetros do experimento (população de 60, 10 espécies, 1000
gerações, 12 episódios por avaliação; os padrões atuais de população e espécies do sistema foram
elevados depois deste experimento)''. A substituição de ``padrão do sistema'' por ``parâmetros do
experimento'' é a correção certa, e o valor de população (60) confere com o artefato que gerou a
Tabela~1 (\texttt{tools/RadarTorres.Experimentos/Resultados/comparacao\_pareada.json},
\texttt{Protocolo.TamanhoPopulacao = 60}). Mas a parte entre parênteses --- ``foram elevados depois
deste experimento'' --- não confere com o histórico git: o commit que elevou população (60→100) e
espécies (6→10), \texttt{c69e4f1}/\texttt{bf5fe32}, é de \textbf{7 de novembro}; o JSON do experimento
traz \texttt{"Data": "2026-11-09 12:56"}, dois dias \textbf{depois}. Os padrões não foram elevados
depois deste experimento: foram elevados antes, e o experimento usou população 60 porque o código que
o gerou (\texttt{Program.cs}, linha 429) fixa esse valor explicitamente em vez de herdar o padrão da
classe --- já então 100 --- enquanto não fixa \texttt{QuantidadeEspecies}, que por isso herdou o
padrão já elevado (10) no momento em que o experimento rodou. É um achado novo, pequeno e sem efeito
sobre a conclusão central, mas é uma inversão factual de cronologia introduzida por esta própria
correção.

\textbf{D.5.} A legenda da Tabela~2 (linhas 637--640) agora diz: ``O banco de dados que originou estes
valores foi posteriormente sobrescrito por execuções seguintes; não há dump congelado para reprodução
independente, e os números aqui são os registrados no momento da observação original.'' Esta auditoria
procurou, nos artefatos versionados do projeto, qualquer arquivo que contivesse os quatro valores da
tabela (9,11; 134,17; 119,62; 137,35) e não encontrou nenhum --- as únicas ocorrências de substrings
parecidas em \texttt{comparacao\_pareada.json} e \texttt{dificuldade\_seeds\_por\_geracao.json} são
coincidências de dígitos dentro de outros números (por exemplo, \texttt{69.117...} e
\texttt{99.111...}), não os valores da Tabela~2. Não há indício de que um dump exista e esteja sendo
omitido. \textbf{Esta é a resposta certa ao achado}: a seção ``O que é achado e o que não é'' deste
mesmo método de auditoria já previa que uma afirmação sem artefato persistido é não auditável, e isso
é limitação da auditoria, não acusação de incorreção --- desde que o texto diga isso com essas
palavras, o que agora faz. D.5 deixa de ser um problema de honestidade e passa a ser apenas uma
limitação declarada de reprodutibilidade, que é precisamente o que a 3ª e a 4ª revisão pediam como
alternativa caso o dump não fosse recuperável.

Os demais achados persistentes da 4ª revisão (D.4, D.6, D.7, D.8, D.9, D.10) foram reconferidos contra
o \texttt{Artigo.tex} e o código em disco hoje e continuam \textbf{idênticos}. Nenhum piorou, nenhum
foi corrigido. A Tabela~1, a tabela central de resultados, não foi tocada por nada desta rodada.

\textbf{Saldo.} O problema de processo que bloqueava a defesa --- duas branches, cada uma com a metade
certa do artigo --- está encerrado, e de forma auditável: um único commit, uma única história. A
resposta a D.5 é o tipo de honestidade que esta banca pede e raramente recebe sem insistência. O preço
de D.2 é pequeno: uma frase que tentou dar contexto histórico ao leitor acertou o quê mas errou o
quando. Nenhum dos três pontos desta rodada muda a conclusão central do trabalho, que continua sendo a
mesma desde a 1ª revisão: o controlador híbrido não superou a heurística pura no cenário testado, e o
texto declara isso no Resumo, no Abstract e na Conclusão.

\section*{B. Nota geral}

\textbf{Nota ponderada: $\approx 8{,}0\,/\,10$} \quad (pareceres anteriores: $6{,}33$; $6{,}20$; $7{,}9$;
$7{,}7$).

\renewcommand{\arraystretch}{1.3}
\begin{longtable}{|p{4.1cm}|p{1.1cm}|p{1.1cm}|p{1.4cm}|p{6.6cm}|}
\hline
\textbf{Critério} & \textbf{Peso} & \textbf{Nota} & \textbf{Contrib.} & \textbf{Justificativa (resumo)} \\
\hline
\endhead
Fundamentação teórica & 20\% & 8,0 & 1,600 &
Inalterada. Nada nesta rodada toca referências ou fundamentação. \\
\hline
Metodologia & 20\% & 8,3 & 1,660 &
\textbf{Sobe de 8,0.} D.2 estava puxando este critério por citar um ``padrão do sistema'' que não
existe mais; a correção resolve isso. Fica abaixo de 8,5 porque a própria correção introduz um erro
de cronologia novo (ver D.2 revisado) e porque D.7 (teto de alvos por quadrante inexistente no
código) persiste. \\
\hline
Resultados e análise & 25\% & 7,8 & 1,950 &
\textbf{Sobe de 7,5.} A Tabela~1 continua reproduzindo exatamente. D.5 deixa de ser uma afirmação
numérica sem fonte e passa a ser uma limitação declarada com honestidade --- not o mesmo que um
artefato recuperado, mas uma melhoria real de rigor. D.6 (recompensa por abate não é fixa em 10)
persiste e segue puxando a nota para baixo. \\
\hline
Redação e estrutura & 15\% & 8,5 & 1,275 &
Inalterada. D.4 (aritmética 7+4 vs.\ 8+3 na Seção de modelo de dados) persiste, idêntica à 4ª
revisão. \\
\hline
Contribuição e originalidade & 10\% & 8,0 & 0,800 &
Inalterada. O resultado negativo permanece declarado no Resumo, no Abstract e na Conclusão. \\
\hline
Reprodutibilidade & 10\% & 7,5 & 0,750 &
\textbf{Sobe de 5,5.} D.1 está resolvido numa única branch auditável, o que é mais do que a 3ª
revisão tinha quando deu 7,0 a este critério (ali a correção só existia em \texttt{origin/main}, sem
confirmação de que alcançaria a branch de entrega). Não chega a 8,5+ porque D.5 continua sem artefato
recuperável (agora declarado, não escondido) e D.9 (execuções fora do repositório) e D.10 (hardware
autodeclarado) persistem sem mudança. \\
\hline
\multicolumn{3}{|r|}{\textbf{TOTAL}} & \textbf{8,04} & \\
\hline
\end{longtable}

\subsection*{Por que a nota voltou a subir}

Porque o mecanismo que a derrubou na 4ª revisão era especificamente processual --- a correção existia,
mas em uma branch que a banca não examinaria. Esse mecanismo desapareceu: há uma única branch agora.
A nota não volta à faixa de 8,3--8,5 que a 4ª revisão projetou porque esta rodada, ao verificar D.2 com
cuidado (comparando datas de commit com a data registrada no JSON do experimento, não apenas lendo a
frase), encontrou um erro novo que a correção introduziu. É um erro pequeno, mas rebaixa levemente
Metodologia em relação ao que seria possível, e por isso o total fica em 8,0 em vez de 8,3+.

\section*{C. Evolução desde o parecer anterior}

\subsection*{C.1 Corrigido ou materialmente melhorado}

\begin{enumerate}
  \item \textbf{D.1 \resolvido} (bifurcação \texttt{main}/\texttt{banco-dados} e, por consequência, a
  correção do genoma versionado e da frase da Seção 2.3.5). \texttt{git merge-base --is-ancestor}
  confirma ancestralidade nos dois sentidos; \texttt{git rev-parse} confirma hash idêntico
  (\texttt{5ef8f91}). A correção que a 1ª, 2ª e 3ª revisões pediram e que a 4ª encontrou isolada na
  branch errada está, agora, na única branch que existe.
  \item \textbf{D.5} (Tabela de certeza estática). A legenda agora declara que o banco de origem foi
  sobrescrito e que não há dump para reprodução independente. Esta auditoria não encontrou nenhum
  artefato versionado com os quatro valores da tabela, o que é consistente com a declaração. É a
  resposta correta a um artefato realmente perdido: declarar o limite, não fabricar o dado.
\end{enumerate}

\subsection*{C.2 Melhorias novas, não pedidas no parecer anterior}

Nenhuma. As três mudanças desta rodada respondem exatamente às três pendências que a 4ª revisão listou
como bloqueantes ou de alto impacto (reconciliação de branches, D.2, D.5); nenhum conteúdo novo além
disso foi encontrado no commit \texttt{5ef8f91}.

\subsection*{C.3 Persiste, idêntico}

\begin{enumerate}
  \item \textbf{D.4} (``sete cadastrais e quatro de treinamento'' vs.\ tabela com 8+3 linhas) ---
  idêntica, linha 524.
  \item \textbf{D.6} (recompensa por abate descrita como 10 pontos fixos; o código soma de 7 a 10) ---
  idêntica, linha 472.
  \item \textbf{D.7} (``cada quadrante admite até cinco alvos simultâneos''; sem teto no código) ---
  idêntica, linha 471.
  \item \textbf{D.8} (``4631 oportunidades de acionamento'' sem indicar que é por execução) ---
  idêntica, linha 609.
  \item \textbf{D.9} (execuções reportadas a partir de cópia em \texttt{C:\textbackslash TEMP}) --- não
  reconferido nesta rodada por ausência de indício de mudança.
  \item \textbf{D.10} (hardware autodeclarado) --- sem mudança.
\end{enumerate}

\subsection*{C.4 Piorou ou surgiu}

\begin{enumerate}
  \item \textbf{D.2 (revisado), cronologia invertida.} A nova frase da linha 235--237 diz que os
  padrões de população e espécies ``foram elevados depois deste experimento''. O commit que os elevou
  (\texttt{c69e4f1}/\texttt{bf5fe32}) é de 7 de novembro; a data registrada no JSON do experimento
  (\texttt{comparacao\_pareada.json}, campo \texttt{Protocolo.Data}) é 9 de novembro. A elevação
  aconteceu \textbf{antes} do experimento, não depois. É um achado novo, introduzido por esta própria
  correção, de gravidade baixa porque não afeta nenhuma conclusão do trabalho --- afeta apenas a
  precisão de uma frase de contexto histórico que o próprio texto optou por incluir.
\end{enumerate}

\section*{D. Problemas}

\renewcommand{\arraystretch}{1.25}
\begin{longtable}{|p{0.9cm}|p{2.6cm}|p{11.6cm}|}
\hline
\textbf{Item} & \textbf{Gravidade} & \textbf{Em uma linha} \\
\hline
\endhead
D.1 & \resolvido & Branches reconciliadas num único commit; genoma fora do versionamento e frase da
Seção 2.3.5 corrigida na branch aberta \\
\hline
D.2 & \baixo\ (revisado) & Valores agora corretos (parâmetros do experimento), mas a frase nova
inverte a cronologia real da elevação dos padrões \\
\hline
D.3 & \retirado & Contagem de tabelas do banco --- argumento de escopo aceito na 3ª revisão; mantido \\
\hline
D.4 & \baixo & ``Sete cadastrais e quatro de treinamento'' não bate com as 8+3 linhas da própria
tabela --- persiste \\
\hline
D.5 & \baixo\ (declarado) & Tabela de certeza estática sem artefato recuperável, agora dito
explicitamente na legenda --- limitação honesta, não mais omissão \\
\hline
D.6 & \moderado & ``10 pontos por alvo interceptado'': o código soma de 7 a 10, conforme a distância
--- persiste \\
\hline
D.7 & \moderado & ``Cada quadrante admite até cinco alvos simultâneos'': não existe esse teto no
código --- persiste \\
\hline
D.8 & \baixo & ``4631 oportunidades de acionamento'' é por execução, não o total --- persiste \\
\hline
D.9 & \baixo & Execuções reportadas a partir de cópia em \texttt{C:\textbackslash TEMP} --- não
reconferido, sem indício de mudança \\
\hline
D.10 & \baixo & Especificações de hardware são autodeclaradas --- persiste \\
\hline
\end{longtable}

\subsection*{D.1 \resolvido\ Branches reconciliadas; correção presente na branch aberta}

\textbf{O que foi anunciado para esta rodada.} Que \texttt{main} e \texttt{banco-dados} foram
reconciliadas por \textit{fast-forward} de \texttt{main} até \texttt{banco-dados}, levando os commits
\texttt{0f58ccb} e \texttt{501ce86} para a branch que está fisicamente aberta.

\textbf{O que esta auditoria confirma.}
\begin{itemize}
  \item \texttt{git merge-base --is-ancestor main banco-dados} retorna sucesso (exit 0).
  \item \texttt{git merge-base --is-ancestor banco-dados main} retorna sucesso (exit 0) --- as duas
  condições juntas só são possíveis se as branches apontam para o mesmo commit.
  \item \texttt{git rev-parse main banco-dados} confirma: as duas resolvem para
  \texttt{5ef8f91c43ed5661e0748e613b2e4d02f3245bf1}.
  \item \texttt{git branch --contains 0f58ccb} e \texttt{git branch --contains 501ce86} listam as duas
  branches.
  \item \texttt{git ls-tree -r HEAD --name-only} filtrado por ``Controlador'' não lista
  \texttt{genoma-controlador.json} nem \texttt{rede-controlador.net} --- apenas código-fonte
  (\texttt{ControladorHibrido.cs}, \texttt{ControladorTreinadoProvider.cs} e afins).
  \item \texttt{git show HEAD:.gitignore} tem a linha
  \texttt{src/RadarTorres.App/Resources/ControladorTreinado/}.
  \item O \texttt{Artigo.tex} em disco, linhas 478--480, tem a frase nova: ``A versão entregue não
  inclui rede treinada; um executável compilado a partir dela usa a heurística pura, comportamento de
  fábrica previsto no projeto.''
  \item O \texttt{.csproj} ganhou um comentário explicativo sobre a condição \texttt{Exists(...)},
  coerente com o estado atual do repositório.
\end{itemize}

\textbf{Uma ressalva, não um achado.} Os dois arquivos (\texttt{genoma-controlador.json},
\texttt{rede-controlador.net}) ainda existem fisicamente no disco desta máquina de auditoria
(\texttt{Test-Path} retorna verdadeiro para os dois). Isso não contradiz a correção: são arquivos
\textit{untracked}, resíduo de uma execução local de treinamento nesta máquina, cobertos pelo
\texttt{.gitignore} e que um \texttt{git clone} limpo não traria. A frase do artigo fala do que o
\textbf{repositório} versiona, não do que pode existir como resíduo local, e isso é exatamente o que o
comentário novo do \texttt{.csproj} explica. Não é achado --- é o comportamento esperado de um
\texttt{.gitignore} bem configurado, registrado aqui para que não seja mal lido como um retorno do
problema original.

\textbf{Por que isto resolve o achado crítico.} O critério desta banca é a afirmação que o exemplar
examinado sustenta. Com uma única branch, não há mais ambiguidade sobre qual exemplar é esse, e a
afirmação ``a versão entregue não inclui rede treinada'' é verdadeira para ele: nenhum arquivo de rede
está rastreado, o \texttt{.gitignore} cobre o diretório, e um clone limpo compila com a heurística pura.

\subsection*{D.2 \baixo\ (revisado) Valores corrigidos, cronologia invertida}

\textbf{O que o texto diz agora.} Linha 235--237: ``parâmetros do experimento (população de 60, 10
espécies, 1000 gerações, 12 episódios por avaliação; os padrões atuais de população e espécies do
sistema foram elevados depois deste experimento)''.

\textbf{O que confere.} O valor de população (60) é o mesmo registrado em
\texttt{tools/RadarTorres.Experimentos/Resultados/comparacao\_pareada.json}
(\texttt{Protocolo.TamanhoPopulacao: 60}), e o código-fonte atual
(\texttt{src/RadarTorres.App/NeuralEvolution/ParametrosTreinamento.cs}, linha 217) tem
\texttt{TamanhoPopulacao = 100} como padrão --- de fato maior que o valor do experimento. A troca de
``padrão do sistema'' por ``parâmetros do experimento'' resolve o problema central que a 3ª e a 4ª
revisão apontaram: o experimento não usa mais nenhum valor que o leitor confundiria com o
comportamento padrão atual do sistema.

\textbf{O que não confere.} A frase entre parênteses afirma que a elevação dos padrões aconteceu
\textit{depois} do experimento. O histórico git mostra o contrário: o commit
\texttt{c69e4f1}/\texttt{bf5fe32} (``Eleva ProbabilidadeAdicionarNo, QuantidadeEspecies (6 -> 10) e
TamanhoPopulacao (60 -> 100)'') tem data de \textbf{7 de novembro de 2026, 01:33}. O JSON do
experimento que produziu a Tabela~1 registra \texttt{"Data": "2026-11-09 12:56"} --- \textbf{dois dias
depois} da elevação dos padrões, não antes. Isto é coerente com o próprio código do experimento
(\texttt{tools/RadarTorres.Experimentos/Program.cs}, linhas 427--431): ele constrói
\texttt{ParametrosTreinamento} fixando \texttt{TamanhoPopulacao = cfg.Populacao} (60, um valor
explicitamente escolhido para o experimento, inferior ao padrão já vigente de 100) mas \textbf{não}
fixa \texttt{QuantidadeEspecies}, que por isso herda o padrão da classe --- já elevado a 10 na data em
que o experimento rodou. Ou seja: a população foi uma escolha deliberada do experimento, diferente do
padrão vigente; as espécies não foram escolhidas, vieram do padrão, que já estava em 10 quando o
experimento rodou e continua em 10 hoje --- não houve elevação alguma de espécies depois deste
experimento, porque a elevação, para as duas variáveis, é a mesma e aconteceu antes.

\textbf{Por que é achado e por que é baixo.} É achado porque é uma afirmação factual sobre o próprio
histórico do projeto, verificável por data de commit e por data de artefato, e a frase afirma o
oposto do que esses dois registros mostram. É baixo porque não muda nenhuma conclusão do trabalho: o
valor de população citado continue correto, e o ponto que a frase tenta fazer --- avisar o leitor que
os valores do experimento não são mais os padrões atuais --- permanece verdadeiro, só que pelo motivo
errado.

\textbf{Correção.} Trocar ``foram elevados depois deste experimento'' por algo como ``o padrão de
população já havia sido elevado para 100 quando este experimento foi executado (7 de novembro) com um
valor menor, escolhido deliberadamente; o padrão de espécies, por sua vez, já estava em 10 e não foi
alterado desde então''. Alternativa mais simples: remover a oração causal e apenas dizer que os valores
citados são os do experimento, não os padrões atuais, sem especular sobre a ordem cronológica.

\subsection*{D.3 \retirado\ Contagem de tabelas do banco de dados}

Mantido retirado, sem reavaliação nesta rodada: o argumento de escopo dos autores, aceito na 3ª
revisão, continua válido.

\subsection*{D.4 \baixo\ Contagem interna inconsistente na mesma frase --- persiste}

Linha 524 do \texttt{Artigo.tex} atual: ``sete cadastrais e de auditoria, e quatro que persistem o
histórico do treinamento'' continua não batendo com a tabela abaixo, que lista 8 linhas de cadastro e
auditoria (Usuários, Preferências do usuário, Câmeras, Objetos detectados, Estados da ocorrência,
Classificações do alvo, Ações realizadas, Modo atual da torre) e 3 de treinamento (Histórico de
gerações, Marcos de \textit{fitness}, Instantâneos do melhor indivíduo). O total (onze) confere; a
partição interna (7+4 no texto vs.\ 8+3 na tabela) não. Idêntico à 3ª e à 4ª revisão.

\subsection*{D.5 \baixo\ (declarado) Tabela de certeza estática sem artefato recuperável}

A legenda (linhas 637--640) agora diz: ``O \textit{fitness} registrado no histórico de treinamento
antes e depois da correção da certeza estática (versão anterior do ambiente). O banco de dados que
originou estes valores foi posteriormente sobrescrito por execuções seguintes; não há dump congelado
para reprodução independente, e os números aqui são os registrados no momento da observação
original.'' Esta auditoria buscou os quatro valores (9,11; 134,17; 119,62; 137,35) em todos os
artefatos JSON versionados do projeto e não encontrou nenhuma ocorrência exata --- apenas coincidências
de substring dentro de números maiores e não relacionados (\texttt{69.117...},
\texttt{99.111...} em \texttt{comparacao\_pareada.json} e \texttt{dificuldade\_seeds\_por\_geracao.json},
que são de outro experimento, sobre outra tabela). Isso é consistente com a declaração de que o banco
de origem foi sobrescrito e não há dump.

Esta é a resposta correta ao achado, pelo próprio critério desta banca: uma afirmação sem artefato
persistido é não auditável, e isso é limitação, não acusação de incorreção --- \textbf{desde que o
texto diga isso}, o que agora faz, com essas palavras. O achado deixa de ser ``número sem fonte
apresentado como se tivesse fonte'' e passa a ser ``limitação de reprodutibilidade declarada'', que é
uma categoria bem mais branda. Mantém-se como baixo, e não como retirado, porque a tabela continua
sendo, objetivamente, a única afirmação numérica central do artigo que nenhum leitor pode verificar de
forma independente --- a correção resolve o problema de honestidade, não o de reprodutibilidade em si,
que é estrutural e sem conserto sem dado novo.

\subsection*{D.6 \moderado\ ``Soma 10 pontos por alvo interceptado'' --- persiste}

Linha 472: a frase ainda descreve a recompensa como fixa. O código
(\texttt{AmbienteTreinamento.cs}, não tocado desde a 4ª revisão) soma
\texttt{RecompensaAbate * (0.7 + 0.3 * bonusDistancia)}, isto é, 7 a 10 pontos. \texttt{git log}
confirma que o último commit a tocar este arquivo é de 17 de novembro, antes de qualquer uma das três
mudanças desta rodada.

\subsection*{D.7 \moderado\ ``Cada quadrante admite até cinco alvos simultâneos'' --- persiste}

Linha 471: frase idêntica. Sem teto por quadrante no código, pela mesma ausência de commits no arquivo
relevante.

\subsection*{D.8 \baixo\ ``4631 oportunidades de acionamento'' --- persiste}

Linha 609: idêntica. Correção inalterada: acrescentar ``em cada execução''.

\subsection*{D.9 \baixo\ Execuções a partir de cópia fora do repositório --- não reconferido}

Sem indício de que os artefatos de resultado tenham sido regenerados nesta janela; mantido como estava.

\subsection*{D.10 \baixo\ Hardware autodeclarado --- persiste}

Sem mudança.

\section*{E. Auditoria de dados e tabelas}

\begin{longtable}{|p{3.6cm}|p{6.6cm}|p{4.4cm}|}
\hline
\textbf{Elemento} & \textbf{O que foi conferido} & \textbf{Veredicto} \\
\hline
\endhead
Relação entre \texttt{main} e \texttt{banco-dados} & \texttt{git merge-base --is-ancestor} nos dois
sentidos, \texttt{git rev-parse} das duas, \texttt{git branch --contains} para \texttt{0f58ccb} e
\texttt{501ce86} & \textbf{Reconciliadas} --- mesmo commit (\texttt{5ef8f91}) \\
\hline
Versionamento do genoma/rede (branch atual) & \texttt{git ls-tree -r HEAD --name-only}, \texttt{git
show HEAD:.gitignore}, grep na linha 478--480 do \texttt{Artigo.tex} em disco, \texttt{Test-Path} dos
dois arquivos & \textbf{Confere} --- não rastreados, \texttt{.gitignore} cobre o caminho, frase nova
presente; arquivos físicos no disco são resíduo local coberto pelo \texttt{.gitignore} \\
\hline
D.2 --- valor de população (60) & \texttt{comparacao\_pareada.json}
(\texttt{Protocolo.TamanhoPopulacao}), \texttt{ParametrosTreinamento.cs} linha 217 & \textbf{Confere
exatamente} \\
\hline
D.2 --- cronologia da elevação dos padrões & \texttt{git log} do commit \texttt{c69e4f1}/\texttt{bf5fe32}
(data 2026-11-07) vs.\ \texttt{Protocolo.Data} do JSON (2026-11-09) & \textbf{Diverge} --- a elevação é
anterior ao experimento, o texto diz o contrário (ver D.2) \\
\hline
D.5 --- Tabela de certeza estática (9,11 / 134,17 / 119,62 / 137,35) & Busca pelos quatro valores em
todos os JSONs versionados do projeto & \textbf{Sem artefato de dados} --- nenhuma ocorrência exata
encontrada, consistente com a declaração de banco sobrescrito \\
\hline
Tabela 1 --- comparação pareada & Não reaberta nesta rodada; nenhum commit tocou o artefato ou o
código que a produz desde a 4ª revisão & \textbf{Sem indício de mudança} --- mantido o veredicto de
confere exatamente \\
\hline
D.4 --- contagem de tabelas (7+4 vs.\ 8+3) & Contagem manual das linhas da Tabela do esquema SQLite no
\texttt{Artigo.tex} (linhas 536--549) & \textbf{Diverge} --- total (11) confere, partição interna não
\\
\hline
D.6 --- recompensa por abate & Não reaberto nesta rodada; \texttt{git log} do arquivo relevante mostra
último commit anterior a todas as mudanças desta rodada & \textbf{Sem indício de mudança} --- persiste
divergência relatada na 3ª/4ª revisão \\
\hline
\end{longtable}

\section*{F. Referências}

Nenhuma referência foi adicionada, removida ou reconferida nesta rodada; nenhum commit desde a 4ª
revisão toca \texttt{referencias.bib} ou as citações no \texttt{Artigo.tex}. \textbf{Não verificadas
nesta rodada}: todas as 21 --- duas foram verificadas na 1ª revisão; as demais seguem não verificadas,
o que não é o mesmo que corretas.

\section*{G. Pontos fortes}

\begin{enumerate}
  \item \textbf{A reconciliação de branches resolveu o problema certo, da forma mais simples
  possível} --- \textit{fast-forward}, sem merge conflituoso, sem reescrita de histórico. É o tipo de
  correção que esta banca pode verificar com dois comandos e confirmar sem ambiguidade.
  \item \textbf{A resposta a D.5 é o padrão correto de honestidade diante de um dado perdido.}
  Declarar a limitação na legenda, em vez de recriar um número ou simplesmente deixar o problema como
  estava, é exatamente o que este método de auditoria pede e nem sempre recebe.
  \item \textbf{A tabela central de resultados continua intocada e reproduzível} --- nenhum achado
  desta rodada a afeta.
  \item \textbf{O resultado negativo permanece declarado no Resumo, Abstract e Conclusão.}
  \item \textbf{Os sete achados abertos restantes são todos de correção textual barata}, nenhum exige
  reexperimento ou dado novo.
\end{enumerate}

\section*{H. Nota final}

\begin{center}
\Large\textbf{8,0 / 10} \quad\normalsize (pareceres anteriores: 6,33; 6,20; 7,9; 7,7)
\end{center}

A nota sobe 0,3 em relação à 4ª revisão porque o mecanismo que a derrubou --- correção presente apenas
numa branch que não era a aberta --- não existe mais: há uma única branch, verificável por dois
comandos de \texttt{git}. Não sobe à faixa de 8,3--8,5 que a 4ª revisão havia projetado porque a
verificação cuidadosa de D.2 (comparando data de commit com data de artefato, não apenas lendo a nova
frase) encontrou um erro de cronologia que a própria correção introduziu. É um erro pequeno e sem
efeito sobre a conclusão central, mas impede a nota plena em Metodologia. Restam seis pendências de
correção textual barata (D.4, D.6--D.10) e uma de cronologia (D.2 revisado); nenhuma exige
reexperimento.

\section*{I. Pendências prioritárias}

\subsection*{Bloqueiam a defesa}

Nenhuma. O único item que bloqueava a defesa na 4ª revisão (branches não reconciliadas) está resolvido.

\subsection*{Alto impacto, esforço baixo}

\begin{enumerate}
  \item Corrigir a frase sobre a cronologia da elevação dos padrões de população/espécies --- dizer
  que a elevação é anterior ao experimento (7 de novembro), não posterior, ou remover a oração causal
  e apenas registrar que os valores citados não são mais os padrões atuais. \emph{(D.2, revisado)}
  \item Trocar ``soma 10 pontos por alvo interceptado'' por ``até 10 pontos (7 a 10, conforme a
  distância)''. \emph{(D.6)}
  \item Remover ``cada quadrante admite até cinco alvos simultâneos'' ou implementar o teto e
  reexecutar. \emph{(D.7)}
\end{enumerate}

\subsection*{Correções de texto}

\begin{enumerate}\setcounter{enumi}{3}
  \item Corrigir a partição ``sete cadastrais e quatro de treinamento'' para bater com as linhas da
  própria tabela (8+3). \emph{(D.4)}
  \item Acrescentar ``em cada execução'' às 4631 oportunidades de acionamento. \emph{(D.8)}
  \item Dizer, na seção de ambiente computacional, que as execuções reportadas partiram de uma cópia
  do projeto fora do repositório. \emph{(D.9)}
  \item Dizer que as especificações de hardware foram coletadas na própria máquina de execução.
  \emph{(D.10)}
\end{enumerate}

\section*{J. Escopo e limitações desta auditoria}

Nenhum arquivo do projeto foi alterado pela banca na produção deste parecer. Todos os comandos usados
foram de leitura: \texttt{git status}, \texttt{git log}, \texttt{git show}, \texttt{git ls-tree},
\texttt{git branch}, \texttt{git merge-base} (com \texttt{--is-ancestor} nos dois sentidos),
\texttt{git rev-parse}, grep no \texttt{.tex} e nos \texttt{.json} em disco, e \texttt{Test-Path} nos
dois arquivos de rede. Nenhum deles escreve no repositório. Nenhuma migração, build ou experimento foi
executado.

Esta rodada \textbf{não} refez a reprodução completa da Tabela 1 nem a varredura de referências feita
nas rodadas anteriores; reconferiu os seis achados persistentes (D.4, D.6--D.10) por grep de linha
contra o texto em disco, não por reexecução de script, apoiada no fato de nenhum commit nesta janela
ter tocado o código ou os dados que os sustentam. Investigou a fundo os três pontos declarados para
esta rodada: reconciliação de branches, D.2 e D.5.

O \texttt{Artigo.log} em disco é de 27/11 14:25, anterior ao commit \texttt{5ef8f91} que contém as três
correções desta rodada; por isso \textbf{não} foi usado para avaliar \texttt{hbox}s ou referências
cruzadas da versão atual, e o artigo não foi recompilado por esta banca para gerar um log novo
(regra de somente leitura).

\textbf{Não verificado nesta rodada:}

\begin{itemize}
  \item \textbf{Recomparação byte a byte do genoma em disco} contra os dez campeões relatados na
  Tabela~1; não é mais relevante para D.1 (o genoma não está mais versionado), mas não foi refeita por
  não ter sido objeto desta rodada.
  \item \textbf{As fontes bibliográficas.} Seguem não verificadas, à exceção das duas confirmadas na
  1ª revisão.
  \item \textbf{Compilação do artigo} e verificação de \texttt{hbox}s/referências cruzadas na versão
  do commit \texttt{5ef8f91} --- o \texttt{.log} disponível é anterior a essa versão.
  \item \textbf{A origem exata dos números da Tabela~1} (confiada ao veredicto ``confere exatamente''
  das rodadas anteriores, não recalculada nesta); o módulo de visão computacional (declarado não
  avaliado pelo próprio artigo); e o firmware do Arduino.
  \item \textbf{O README da branch \texttt{genes-treinamento}}, não reaberto por ausência de commit
  novo nela.
\end{itemize}

\textbf{Artefatos consultados nesta rodada:} \texttt{Artigo/Texto/Artigo.tex} (em disco, HEAD
\texttt{5ef8f91}); histórico git completo (\texttt{git log}, \texttt{--oneline}, com datas;
\texttt{git log --follow -p} para \texttt{ParametrosTreinamento.cs}); \texttt{git show} dos commits
\texttt{0f58ccb}, \texttt{501ce86}, \texttt{c69e4f1}/\texttt{bf5fe32}; \texttt{git ls-tree -r HEAD};
\texttt{.gitignore}; \texttt{src/RadarTorres.App/RadarTorres.App.csproj};
\texttt{src/RadarTorres.App/NeuralEvolution/ParametrosTreinamento.cs};
\texttt{tools/RadarTorres.Experimentos/Program.cs};
\texttt{tools/RadarTorres.Experimentos/Resultados/comparacao\_pareada.json};
\texttt{tools/RadarTorres.Experimentos/Resultados/dificuldade\_seeds\_por\_geracao.json};
\texttt{src/RadarTorres.App/Resources/ControladorTreinado/\{rede-controlador.net,
genoma-controlador.json\}} (existência em disco); \texttt{Software/Banco\_de\_Dados/} (busca por
dump/db versionado); \texttt{Artigo/Revisoes/Parecer\_Banca\_TCC\_rev4.tex} (parecer anterior, íntegro).

\end{document}
